% Chapter 1

\chapter{Introducción general} % Main chapter title

\label{Chapter1} % For referencing the chapter elsewhere, use \ref{Chapter1} 
\label{IntroGeneral}

En la actualidad, la banca digital se fundamenta en un paradigma de interacción táctil y visual que, si bien es eficiente para la mayoría de los usuarios, impone barreras significativas para personas con discapacidades visuales o motoras.
La solución desarrollada en este proyecto está orientada a los usuarios de la banca minorista del banco Galicia, buscando empoderar a aquellos con diversidades funcionales y ofrecer una experiencia de usuario más ágil y con menor fricción operativa.
El valor central de esta solución reside en la convergencia de la inclusión financiera con la máxima seguridad y privacidad de la información. A diferencia de los asistentes virtuales tradicionales que dependen del procesamiento en la nube y exponen datos sensibles, este sistema implementa un pipeline de Inteligencia Artificial de ejecución estrictamente local. Esto garantiza que tanto los datos acústicos como las intenciones transaccionales nunca abandonen el dispositivo del usuario.

%----------------------------------------------------------------------------------------

% Define some commands to keep the formatting separated from the content 
\newcommand{\keyword}[1]{\textbf{#1}}
\newcommand{\tabhead}[1]{\textbf{#1}}
\newcommand{\code}[1]{\texttt{#1}}
\newcommand{\file}[1]{\texttt{\bfseries#1}}
\newcommand{\option}[1]{\texttt{\itshape#1}}
\newcommand{\grados}{$^{\circ}$}

%----------------------------------------------------------------------------------------

\section{Contexto de la banca móvil{}}

En la última década, la banca digital ha experimentado una transformación radical, migrando de plataformas web tradicionales a aplicaciones móviles que se han consolidado como el principal punto de contacto entre los clientes y las instituciones financieras. En este entorno de alta competitividad, la innovación continua en la experiencia de usuario se ha vuelto un diferenciador estratégico clave para la retención y satisfacción del cliente. Las interfaces conversacionales, impulsadas por los recientes avances en Inteligencia Artificial y Procesamiento de Lenguaje Natural, representan la siguiente frontera en la interacción humano-computadora, permitiendo una comunicación más fluida, rápida e intuitiva. Este proyecto se enmarca en dicha evolución tecnológica, explorando la viabilidad de aplicar tecnologías de reconocimiento y comprensión de voz para la ejecución de operaciones financieras directamente desde el ecosistema de una aplicación de banca móvil.


%----------------------------------------------------------------------------------------

\section{Estado del arte}

En la actualidad, los asistentes de voz de propósito general han alcanzado un grado de madurez notable en tareas de reconocimiento automático de habla y comprensión del lenguaje natural. Sin embargo, la adopción e integración de estas tecnologías en dominios de misión crítica y alta seguridad, como el sector bancario, sigue siendo limitada. Las soluciones bancarias que han incorporado interfaces conversacionales suelen basarse en integraciones en la nube, donde los flujos de audio y datos se envían a servidores externos para su procesamiento. Este enfoque tradicional genera legítimas preocupaciones respecto a la privacidad, la latencia y la exposición de información financiera sensible.
Para mitigar estos riesgos, la tendencia tecnológica emergente apunta hacia el procesamiento en el dispositivo. Bajo este paradigma, los modelos de Inteligencia Artificial se ejecutan localmente utilizando el hardware del dispositivo móvil, garantizando que los datos acústicos y las intenciones transaccionales del usuario nunca abandonen su teléfono. Esta arquitectura maximiza la seguridad y la confidencialidad, pilares fundamentales para la viabilidad de cualquier solución dentro de la industria financiera.



%----------------------------------------------------------------------------------------

\section{Objetivos y alcance}

\subsection{Objetivo general}

El objetivo principal de este trabajo es diseñar, desarrollar y evaluar un prototipo funcional de un asistente transaccional por voz integrado en el entorno de una aplicación móvil para iOS. Este prototipo servirá como una prueba de concepto destinada a validar la viabilidad técnica de un sistema que permita a los usuarios ejecutar operaciones financieras mediante comandos de voz en lenguaje natural.

\subsection{Objetivos específicos}

Para alcanzar el objetivo general, se plantean los siguientes objetivos específicos:

\begin{itemize}
\item \textbf{Mejorar la accesibilidad e inclusión} 

Desarrollar un canal de interacción alternativo que empodere a personas con discapacidades visuales o motoras, reduciendo la dependencia de la navegación táctil tradicional.
\item \textbf{Garantizar la seguridad y privacidad de los datos} 

Implementar modelos de Machine Learning optimizados para ejecutarse de manera estrictamente local mediante Core ML de Apple, asegurando que el audio y los datos sensibles del usuario se procesen íntegramente en el dispositivo. 
\item \textbf{Optimizar la experiencia de usuario}

Simplificar operaciones bancarias complejas condensándolas en un único paso conversacional, reduciendo de esta manera la fricción y el tiempo necesario para completar una transacción.
\item \textbf{Sentar bases técnicas para futuras auditorías}

Generar la documentación arquitectónica y técnica necesaria para facilitar una futura evaluación y validación por parte del equipo de Ciberseguridad de la entidad bancaria.
\end{itemize}

\subsection{Alcance y limitaciones}

Con el propósito de mantener el enfoque en la validación tecnológica del motor de Inteligencia Artificial y la experiencia de usuario, el proyecto se delimita bajo parámetros estrictos de alcance. 


\begin{itemize}
\item \textbf{Dentro del alcance:} 

El desarrollo contempla exclusivamente la creación de una aplicación prototipo nativa para iOS, la cual operará de manera independiente y aislada. El núcleo técnico abarcará la implementación del pipeline de Inteligencia Artificial en el dispositivo físico, integrando la captura segura de audio, un modelo de reconocimiento automático de habla y un módulo de interpretación de lenguaje natural para extraer intenciones y entidades clave. Dado el carácter de la prueba de concepto, todo el flujo de transacción, la pantalla de confirmación y los resultados de las operaciones se simularán utilizando datos completamente ficticios.
\item \textbf{Fuera del alcance:} 

El proyecto no contempla la integración del prototipo, ni de sus componentes, con el código fuente de la aplicación productiva oficial, ni la conexión con APIs, servicios de backend o bases de datos reales de la institución financiera. Asimismo, queda excluida la puesta en marcha en producción y la publicación en la App Store. Desde el punto de vista del manejo de la información, está estrictamente excluido el uso o acceso a datos reales de clientes, así como la ejecución de transacciones que movilicen dinero real. Finalmente, el sistema se limitará al español rioplatense focalizado, dejando fuera del alcance el soporte multilingüe y operaciones financieras distintas a transferencias y consultas.


\end{itemize}


\section{Motivación: accesibilidad e inclusión}
El paradigma de diseño predominante en las aplicaciones de banca móvil actuales se basa casi exclusivamente en la interacción táctil y visual. Si bien este enfoque resulta eficiente para la mayoría de los usuarios, presenta barreras tecnológicas y de usabilidad significativas para grupos demográficos importantes. Por un lado, las personas con discapacidades o disminuciones visuales enfrentan el desafío de navegar a través de menús complejos, campos de texto y botones, tareas que a menudo resultan frustrantes incluso con el soporte de las herramientas de accesibilidad nativas de los sistemas operativos. Por otro lado, las personas con discapacidades motoras encuentran que la necesidad de una manipulación precisa de la pantalla dificulta o impide directamente el uso autónomo de la aplicación para realizar operaciones financieras cotidianas.
La motivación principal de este trabajo es derribar estas barreras tecnológicas, promoviendo una inclusión financiera real mediante el desarrollo de un canal de interacción alternativo. La implementación de un asistente transaccional por voz no solo genera un impacto social positivo al mejorar radicalmente la accesibilidad, sino que también ofrece una experiencia de usuario superior, simplificando operaciones complejas en un único paso conversacional y reduciendo la fricción del proceso.
